\section{Analysis}
\label{sec:analysis}

We organize the analysis in the same causal order as the method: support coverage, teacher-selected adjacency, multi-hop targets, cross-anchor consistency, and candidate efficiency. Aggregate downstream changes are interpreted together with full-corpus relational diagnostics; they are not used alone as evidence for a mechanism.

\subsection{Batch-Selected Teacher-Mass Exposure}

To test Proposition~\ref{prop:batch-exposure}, we track exposed teacher mass against relation opportunities for batch- and teacher-selected support. During batch-local training, we compare in-batch KL with full-corpus JS divergence and teacher-neighbor recall. Coverage-only probes at multiple corpus sizes test the predicted dependence on corpus scale without confounding it with student optimization.

The relevant empirical question is whether teacher-selected candidates capture substantially more mass at the same number of scored relations. If an empirical teacher target is diffuse, the retained-mass curve---rather than the concentrated-support special case---determines the strength of this advantage.

\subsection{Matched-Compute Support Selection}

The core comparison uses batch-local KD, random corpus support, RIPPLE one-step support, RIPPLE with diffusion, full RIPPLE with row supervision, and a degree-preserving random graph. Batch $\rightarrow$ random corpus tests broader exposure; random corpus $\rightarrow$ RIPPLE one-step tests teacher relevance at a matched number of anchor--column scores. The remaining variants isolate diffusion, row coupling, and generic graph regularization. All variants share the student initialization, optimizer, schedule, seeds, and scored-pair budget; Appendix~\ref{app:ablation-protocol} specifies the controls.

\begin{figure}[t]
    \centering
    \IfFileExists{figures/figure2_core_ablation.pdf}{%
        \includegraphics[width=\linewidth]{figures/figure2_core_ablation.pdf}%
    }{%
        \fbox{\parbox[c][0.18\textheight][c]{0.95\linewidth}{\centering
        \textbf{Matched-compute analysis.}\\[2mm]
        Batch-local, random corpus, RIPPLE one-step, RIPPLE diffusion,\\
        full RIPPLE, and a degree-preserving random graph.\\[1mm]
        Downstream Avg.$\uparrow$ \quad Full-corpus JS$\downarrow$ \quad Held-out $E_8\downarrow$.}}
    }
    \caption{\textbf{Matched-compute analysis of RIPPLE.} Random corpus support controls for broader exposure and the number of comparisons; one-step RIPPLE isolates teacher-selected adjacency; diffusion and row-supervision variants test the remaining components; the random graph controls for generic graph regularization. Bars report mean $\pm$ standard deviation over seeds.}
    \label{fig:core-ablation}
\end{figure}

\subsection{Multi-Hop and Cross-Anchor Transfer}

Diffusion should improve relations beyond the directly supervised one-hop band rather than merely increase neighborhood coverage. We therefore evaluate held-out diffusion errors at broader scales together with teacher--student JS in fixed teacher-rank bands. Varying only graph width tests whether additional teacher support improves these fixed bands.

For row supervision, we compare diffusion-only training with full RIPPLE and measure reciprocal-overlap disagreement alongside full-corpus JS. A downstream improvement without reduced disagreement would not support the proposed cross-anchor mechanism.

\subsection{Coverage, Calibration, and Efficiency}

We vary the candidate budget and measure retained teacher mass, selected-versus-unselected mass calibration, student off-support mass, sampled-to-full gradient agreement, downstream quality, step time, and peak memory. To test Proposition~\ref{prop:global-geometry}, we additionally track selected-support cosine distortion $\epsilon_T$, missed teacher mass $\bar\delta_T$, their bound $\epsilon_T+\bar\delta_T$, and full-corpus distortion $\mathcal E_T$ over training epochs. High retained mass with small restricted KL is insufficient if mass calibration remains poor; these quantities are therefore reported separately. Saturation at a small candidate budget before cost approaches full scoring is the evidence required for concentration-adaptive computation.

\begin{figure}[t]
    \centering
    \IfFileExists{figures/figure3_mechanisms.pdf}{%
        \includegraphics[width=\linewidth]{figures/figure3_mechanisms.pdf}%
    }{%
        \fbox{\parbox[c][0.16\textheight][c]{0.95\linewidth}{\centering
        \textbf{Mechanism and efficiency diagnostics.}\\[2mm]
        (a) graph width vs. retained mass and fixed rank-band JS.\\
        (b) candidate budget vs. calibration, gradient fidelity, time, and memory.}}
    }
    \caption{\textbf{Coverage, calibration, and efficiency.} Left: graph width is evaluated on fixed teacher-rank bands to distinguish broader transfer from trivially increased support. Right: candidate budget is evaluated through target coverage, mass calibration, gradient fidelity, downstream quality, and computational cost.}
    \label{fig:mechanisms}
\end{figure}

The controlled settings are summarized in Appendix~\ref{app:diagnostic-protocols}. Until these mechanism diagnostics are populated, the main table supports effectiveness but not a causal attribution to relation selection, diffusion, or row coupling.
